\section{模型改进与推广}

\subsection{蒸发潜热对照（M1，条件性扩展）}
\label{sec:M1}

问题2/3 的 M0 基线\textbf{不含蒸发潜热}。作为对照，M1 在表面能量平衡中加入潜热项
\begin{equation}
  -k\frac{\partial T}{\partial r}\bigg|_{R}=h\left(T_s-T_\infty\right)+\lambda J_w,
  \qquad
  J_w=\rho_d\,h_m\left(C_s-\Cinf\right),
\end{equation}
其中 $J_w$ 的单位是 kg/(m$^2\cdot$s)，\textbf{不得}把 $h_m(C_s-\Cinf)$
直接乘潜热当热通量——它的单位是 kg/(kg$\cdot$s)，不是 kg/(m$^2\cdot$s)，
必须乘密度基准 $\rho_d$。

\begin{table}[H]
  \centering
  \caption{M0 与 M1 的对照}
  \label{tab:M1}
  \footnotesize
  \begin{tabular}{@{}lccc@{}}
    \toprule
    项 & M0 & M1 & 差 \\
    \midrule
    3 h 中心温度 & $49.849\ ^\circ$C & $31.786\ ^\circ$C & $\mathbf{-18.06}$ K \\
    3 h 表面温度 & $49.967\ ^\circ$C & $32.402\ ^\circ$C & $-17.56$ K \\
    3 h 中心含水率 & 1.7662 & 2.2515 & $+0.485$（\textbf{更慢}） \\
    $\tstar$ & 57.49 h & 60.33 h & $\mathbf{+2.84}$ h（$\times1.049$） \\
    \bottomrule
  \end{tabular}
\end{table}

\begin{quote}
  \textbf{与文献的定量交叉验证}：文献 [1b] 报道``忽略潜热会使初期物料温度预测
  偏高 $12\sim22\ ^\circ\mathrm{C}$''。本项目独立复现的差值是 $\mathbf{18.06}$ K，
  \textbf{落在该区间内}——一次独立定量印证。
\end{quote}

\begin{figure}[H]
  \centering
  \includegraphics[width=0.75\textwidth]{FIG-43_latent.png}
  \caption{潜热开关对照}
  \label{fig:FIG43}
\end{figure}

\textbf{为什么不能据此修订问题3 答案}：M1 的量级\textbf{正比于干基密度 $\rho_d$}，
而 $\rho_d$ 在本题中\textbf{没有给定}：

\begin{itemize}[leftmargin=1.6em, itemsep=2pt]
  \item 附录3 的 $\rho(C)=650+128C$ 是\textbf{湿物料体积密度}；
  \item 由它反推的 $\rho_d=\rho(C)/(1+C)$ 在 $C=0$ 给 650、在 $C=2.55$ 给 275.0
        ——\textbf{不是常数}，与``干物质守恒''\textbf{不兼容}；
  \item 本项目取\textbf{初始状态基准} $\rho_d=976.4/3.55=275.042$ kg/m$^3$
        并\textbf{显式声明}。
\end{itemize}

$\rho_d$ 翻倍则潜热效应翻倍（严格正比，相邻比 $2.03\sim2.07$）：

\begin{table}[H]
  \centering
  \footnotesize
  \begin{tabular}{@{}lccc@{}}
    \toprule
    $\rho_d$ 倍率 & $\rho_d$ / (kg/m$^3$) & $\tstar$ / h & $\Delta$ vs M0 \\
    \midrule
    0.25 & 68.76 & 58.17 & $+0.68$ h \\
    0.50 & 137.52 & 58.87 & $+1.38$ h \\
    \textbf{1.00（本文采用）} & \textbf{275.04} & \textbf{60.33} & $\mathbf{+2.84}$ h \\
    2.00 & 550.08 & 63.37 & $+5.88$ h \\
    \bottomrule
  \end{tabular}
\end{table}

\begin{quote}
  \textbf{结论}：M1 只能写成``\textbf{在本文声明的 $\rho_d$ 取值下}，
  蒸发潜热使 3 h 中心温度低约 18 K、使 $\tstar$ 延长约 2.8 h''。
  \textbf{不得}写成``药材真实烘干时间应为 60.3 h''，
  也\textbf{不得}用 M1 替换问题3 的答案。
\end{quote}

\subsection{烘房系统模型（M2 工程推广）}

本文把烘房环境作为\textbf{给定边界}（单向驱动），这是 M0 的合理简化。
若需预测新工况（改变装料量、换气量、加热策略），
应升级为\textbf{烘房—物料双向耦合系统模型}。正确的水分收支形式
（简化充分混合空气）为
\begin{equation}
  \frac{\mathrm d}{\mathrm dt}\left(M_{da}Y\right)
  =\dot m_{da,\mathrm{in}}Y_{\mathrm{in}}
   -\dot m_{da,\mathrm{out}}Y+\sum_i A_i J_{w,i}-\dot m_{\mathrm{cond}},
\end{equation}
还需\textbf{干空气质量平衡}。

\begin{quote}
  若 $M_{da}$ 不变，才可写成 $M_{da}\dot Y$；
  \textbf{不能用 $V_{air}\dot Y$ 等于质量流率的形式}——量纲不闭合。
\end{quote}

\textbf{数据需求}：空气干质量、装料量、加热与排湿条件、
新风/排风干空气质量流量或除湿速率。
\textbf{验证方式}：选择实测环境作为验证数据，
\textbf{不能同时将环境固定与自由求解}。

\subsection{品质后处理（M2 推广）}

有目标成分动力学后，可沿材料轨迹计算成分保留分数：
\begin{equation}
  \frac{q(t)}{q(0)}=\exp\left[-\int_0^t k_q\!\left(T(\tau),C(\tau)\right)\mathrm d\tau\right].
\end{equation}
这是\textbf{待标定的候选动力学}，不是文献为本题药材提供的通用公式。

\begin{table}[H]
  \centering
  \footnotesize
  \begin{tabular}{@{}ll@{}}
    \toprule
    扩展方向 & 额外数据需求 \\
    \midrule
    微生物 & 水活度关系 $a_w(C,T)$、菌种、初始菌量 \\
    结构（结壳 / 开裂） & 应变、本构关系、强度参数 \\
    挥发成分保留 & 成分种类、组织结构、动力学参数 \\
    \bottomrule
  \end{tabular}
\end{table}

\textbf{无水活度关系时}：只能说明``含水率达标 $\ne$ 微生物合格''，
\textbf{不能制定通用杀菌温度}。

\subsection{结构力学（M2 路线）}

建模链路（仅作路线描述）：
湿度场 $\to$ 收缩应变 $\to$ 本构关系 $\to$ 应力场 $\to$ 断裂判据 $\to$ 诊断量。

\begin{quote}
  不能由湿度梯度直接算裂纹概率或裂纹位置，
  也不能把``多缩 $7\%$''当作仿射收缩的验证。
  可输出梯度作为\textbf{非标定诊断量}，但不叫``裂纹概率''。
  \textbf{题目半径数据不证明内部均匀收缩。}
\end{quote}
